\documentclass[10pt,a4paper]{amsart}
\usepackage[margin=19mm]{geometry}
\usepackage{amsmath,amssymb,mathtools}
\usepackage[scr=rsfs,cal=euler]{mathalfa}
\usepackage{xfrac}
\usepackage[unicode,hidelinks,bookmarks=false]{hyperref}
\newcommand{\ie}{i.e.\ }
\newcommand{\cf}{cf.\ }
\newcommand{\resp}{respectively\ }
\newcommand{\Subsection}{\subsection}
\providecommand{\nobreakdash}{\nobreak\hbox{-}}
\setlength{\emergencystretch}{2em}
\multlinegap0pt
\usepackage{xcolor}
\usepackage[shortlabels]{enumitem}
\usepackage[normalem]{ulem}
\newtheorem{prop}{Proposition}
\newtheorem{lemma}{Lemma}
\newtheorem{pf}{Pf}
\def\Id{\rm{Id}}
\def\Re{{\textrm{Re}}}
\title{On the rate of convergence to steady state\\ in a linear chromatography model}
\author{Joaqu\'in Menacho, Marta Pellicer, J. Sol\`a-Morales}
\date{Source: arXiv:2605.15973v1 (CC BY 4.0)}
\begin{document}
\maketitle

\noindent\emph{Source and licence.} On the rate of convergence to steady state\\ in a linear chromatography model. Version arXiv:2605.15973v1 (CC BY 4.0), distributed by arXiv under the Creative Commons Attribution 4.0 International licence, http://creativecommons.org/licenses/by/4.0/, as recorded in the arXiv licence metadata for this exact version and shown on the abstract page https://arxiv.org/abs/2605.15973v1. Full licence notice: \texttt{licenses/arxiv-2605.15973-CC-BY-4.0.txt}.\par\smallskip

\noindent\textit{Editorial scope.} Counted unit: the article's lemma invers (source0.tex lines 600--602). The article's complete original proof of it is delivered with it (source0.tex lines 603--640, one \texttt{proof} environment of 38 lines). No mathematics is written, repaired or re-derived: the statement and the proof are the article's own lines, in the article's own notation, and no retained line is substituted or re-flowed.  Retained uncounted as the source-local context the proof needs: lines 186--188; lines 220--227; lines 306--310; lines 333--335; lines 381--383; lines 401--401; lines 447--451.  Nothing was omitted from any retained span, so the package carries no omission record.  No retained line contains a citation, so the wrapper carries no bibliography and no bibliography entry.  No retained line contains a figure environment, an image-inclusion command or drawing code, so no image is delivered and the package declares no asset dependency.  Editorial formatting only, in addition: an AMS \texttt{amsart} preamble that restates the article's own declarations and macros used by the retained lines, namely the environments \texttt{prop}, \texttt{lemma}, \texttt{pf} on separate counters, together with the standard packages those lines need.  The retained environments print in this document as Proposition 1, Lemma 1, Pf 1 in order of appearance, whereas the article prints the same items with the numbering of its own full document.  The complete native source of the article as distributed in the arXiv e-print for this exact version is bundled unchanged as \texttt{sources/200695/source0.tex}.
\begin{equation}\label{ineq}
v_1>v_2>0, \ v_2<v_3, \ v_3>v_4>0 \text{ and } v_4<v_1.
\end{equation}

\begin{equation}\label{bcn0}
\left\{ \begin{array}{lll}
  &c(x_i^-)=c(x_i^+) & \hbox{for}\: x_i=-1,1,\\
  &v_1c(-2^+)=v_4c(2^-),&\\
  &v_3c(0^+)= v_2c(0^-),& \\
  &q(x_i^-)=q(x_i^+) & \hbox{for}\: x_i=-1,0,1,\pm 2 \hbox{. ($2^+$ means $-2^-$)}.
\end{array}\right.
\end{equation}

\begin{equation} \label{opA}
\left.\left(A\begin{pmatrix}c\\q\end{pmatrix}\right)\right|_{I_i}=A_i \left(\left.\begin{pmatrix}c\\q\end{pmatrix}\right|_{I_i}\right):=\begin{pmatrix}
-v_i\frac{d}{dx}-R P^2 \,\Id& RP\, \Id\\RP\, \Id&\frac{d}{dx}-R\, \Id
\end{pmatrix}\begin{pmatrix}c\\q\end{pmatrix}\ \textrm{for }  I_i=[x_i,x_{i+1}],\ i = 1,2,3,4.
\end{equation}

\begin{equation}\label{eq:maxmon}
{\rm Re}(\langle A(c,q)^T, (c,q)^T\rangle_X\le 0
\end{equation}

\begin{prop}\label{prop:Reneg}
Let $\lambda$ be any eigenvalue of $A$ defined in \eqref{opA}. If inequalities \eqref{ineq} hold (even if they are not strict), then $\Re(\lambda)\leq 0$. Moreover, if at least one of the four dissipative inequalities \eqref{ineq} holds in the strict sense, then $\Re(\lambda)<0$.
\end{prop}

\section{The Characteristic Equation}\label{sec:carac}

\begin{equation}\label{eq:C}
C(\lambda):= M_1(\lambda,-1^-)\cdot
\begin{pmatrix} v_4/v_1&0\\0&1\end{pmatrix}\cdot M_4 (\lambda, 2^-)\cdot M_3(\lambda, 1^-)\cdot
\begin{pmatrix} v_2/v_3&0\\0&1\end{pmatrix}\cdot M_2 (\lambda,0^-),
\end{equation}

\begin{lemma}\label{invers}
For $s\in(0,+\infty)$ the linear operator $(-A+s \Id)^{-1}$ exists, and is a bounded and compact operator from $X$ to $X$ with $\|(-A+s \Id)^{-1}\|\le s^{-1}$.
\end{lemma}

\begin{pf}
Given $(\varphi_1,\varphi_2)^T\in X$ we have to find $(c,q)^T\in D(A)$ such that $-A(c,q)^T+s(c,q)^T=(\varphi_1,\varphi_2)^T$, that is
\begin{equation}\label{eqn1nh}
\dfrac{d}{dx}
\begin{pmatrix}
c\\q
\end{pmatrix}
=
\begin{pmatrix}
-\dfrac{s+P^2R}{v_i}&\dfrac{RP}{v_i}\\-RP&s+R
\end{pmatrix}
\begin{pmatrix}
c\\q
\end{pmatrix}+\begin{pmatrix}-\dfrac{1}{v_i}\varphi_1\\ \varphi_2\end{pmatrix} \hspace{0.25cm} \textrm{      in each interval $I_i=[x_i,x_{i+1}]$} .
\end{equation}
We start with the case when the functions $\varphi_1,\varphi_2$ are continuous on each closed interval $I_i$. In this case, this is a standard non-homogeneous linear system of ordinary differential equations on each interval $I_i$. As in Section \ref{sec:carac}, the simplest way to solve it with the boundary conditions \eqref{bcn0} is perhaps in the order $I_2,I_3,I_4, I_1$ (that is, starting at $x=-1^+$). Observe that the initial condition of the equation in the next interval is obtained after the final value in the previous interval (and imposing the boundary conditions \eqref{bcn0}). Then, at the end, one has to merely impose that $(c(-1^-),q(-1^-))^T$ coincides with $(c(-1^+),q(-1^+))^T$.

To show that such a solution exists, we can construct a solution of the form $(c,q)^T = (c^h,q^h)^T + (c^p,q^p)^T  $, where $(c^h,q^h)^T$ is a solution of the homogeneous version of system \eqref{eqn1nh}, and $(c^p,q^p)^T$ is a particular solution of system \eqref{eqn1nh}. In both cases, we impose that the boundary conditions \eqref{bcn0} must be fulfilled, but we still have to impose the conditions at the first end $x=-1^+$.

In the case of the particular solution, we impose $(c^p(-1^+),q^p(-1^+))^T=(0,0)^T$, for example. Then, the particular solution $(c^p,q^p)^T$ exists in each interval $I_i$ because of the Existence and Uniqueness Theorem. In the case of the solution of the homogeneous system, we impose $(c^h(-1^+),q^h(-1^+))^T=(c_1,q_1)^T$, where $(c_1,q_1)^T$ is the algebraic solution of the linear system
$$C(s)(c_1,q_1)^T+(c^p(-1^-),q^p(-1^-))^T=(c_1,q_1)^T$$
with $C(s)$ being the Return Map given by the $2\times2$ matrix defined in \eqref{eq:C}. Recall that applying $C(s)$ to any solution at $x=-1^+$ means to run a whole cycle, that is,  $C(s)(c_1,q_1)^T =C(s)(c^h(-1^+),q^h(-1^+))^T = (c^h(-1^-),q^h(-1^-))^T$. We can assure that the above linear system has an algebraic solution since we know that the matrix $C(s)$ has not the eigenvalue $1$ when $s>0$ (because, in this case, $\lambda=s>0$ would be an eigenvalue of $A$, which is not possible, see Proposition \ref{prop:Reneg}). Hence, again by the Existence and Uniqueness Theorem for ODEs, $(c^h,q^h)^T$ exists. Observe that, with the previous definitions, we have $(c(-1^+),q(-1^+))^T = (c(-1^-),q(-1^-))^T$, and $(c,q)^T$ is a solution of the system \eqref{eqn1nh} with boundary conditions \eqref{bcn0}, as desired.

Still with the functions $\varphi_1,\varphi_2$ being continuous on each $I_i$,  we multiply the equality $-A(c,q)^T+s(c,q)^T=(\varphi_1,\varphi_2)^T$ by $(c,q)^T$ with the scalar product $\langle\cdot,\cdot\rangle_X$, and we get that
$$\langle-A(c,q)^T,(c,q)^T\rangle_X+s\langle (c,q)^T,(c,q)^T\rangle_X=
\langle (\varphi_1,\varphi_2)^T,(c,q)^T\rangle_X.$$
Then, we can take the real part of the previous equality and recall that $\Re\left(\langle-A(c,q)^T,(c,q)^T\rangle_X\right)\ge 0$ (see \eqref{eq:maxmon}). Consequently, by using the Cauchy-Schwarz inequality, we obtain
$$
s\|(c,q)^T\|_X^2\le \Re\left(\langle (\varphi_1,\varphi_2)^T,(c,q)^T\rangle_X\right)\le\|(\varphi_1,\varphi_2)^T\|_X\|(c,q)^T\|_X.
$$
From this we get the useful estimate $\|(c,q)^T\|_X\le s^{-1}\|(\varphi_1,\varphi_2)^T\|_X$,
%\begin{equation*}%\label{est}
%\|(c,q)^T\|_X\le s^{-1}\|(\varphi_1,\varphi_2)^T\|_X
%\end{equation*}
that means the continuity of $(-A+s\Id)^{-1}$ from $X$ to $X$ when restricted to continuous functions in each $I_i$.

Since these continuous functions are dense in $X$, the operator $(-A+s\Id)^{-1}$ can be extended to a continuous operator from $X\to X$. By reasoning along sequences $(\varphi_1^{k},\varphi_2^{k})$ which are convergent in $X$, we see that the extended operator $(-A+s\Id)^{-1}$ gives precisely the solutions of \eqref{eqn1nh} when the derivatives ${d}/{dx}$ are understood in the usual weak-$H^1$ sense. This in particular implies that $(-A+s\Id)^{-1}(X)\subset D(A)$. One also easily sees that  $ (-A+s \Id)^{-1}$ is bounded from $X$ to $D(A)$, when $D(A)$ is equipped with the piecewise $H^1$ norm. And then the compactness from $X$ to $X$ follows from the compactness of the embedding $D(A)\subset X$.
\end{pf}

\end{document}
